% ===============================================================
%  Chapter 10 -- Problem Statement                       (Part II)
% ===============================================================
\chapter{Problem Statement: The Cost of Forced Writebacks on Dirty Sharing}
\label{ch:problem}

\begin{figure}[htbp]
\centering
\adjustbox{max width=\linewidth}{%
\begin{forest}
  chapter map,
  [Chapter 10
    [The problem [In words] [Formally]]
    [Scope {\&} assumptions]
    [Research questions [RQ1 to RQ4]]
    [Requirements for a solution [Decision tree]]
  ]
\end{forest}}
\caption[]{Structure of Chapter~10: the problem in words
and formally, its scope, and the research questions Parts~III to~VI answer.}
\label{fig:ch10-map}
\end{figure}

% ---------------------------------------------------------------
\section{The Problem}
\label{sec:prob-sentence}
\label{sec:prob-formal}
\sectionmap{[The problem [MESI writes back on every dirty sharing transfer]
  [Not required by the axioms] [{Measured by B, T, AMAT and F}]]}
\index{MESI}\index{forced writeback}\index{dirty sharing}\index{write-back cache}\index{AMAT}\index{snooping protocol}\index{directory protocol}\index{sharing patterns!producer--consumer}\index{coherence axioms}%

\begin{problembox}[Problem Statement]
In a multicore processor that uses MESI, whenever one core reads a block that
another core holds modified, the protocol writes the block back to memory as well
as supplying it to the reader. The coherence axioms do not require this writeback;
MESI's rule that shared blocks are clean does. The writeback uses interconnect and
memory bandwidth, is often overwritten before memory's copy is read, and its share
of total traffic varies widely with the sharing pattern. We ask how much the
forced writeback costs on 2 to 16 cores, how much of that cost a coherent protocol
can remove, and whether the benefit of removing it can be predicted from a
workload's behaviour under MESI alone.
\end{problembox}

The system is a chip with $N\in\{2,4,8,16\}$ cores, each with a private
write-back L1 cache, sharing an LLC and main memory over a snooping bus or a
directory network, with 64-byte blocks. The workload $W$ is a shared-memory
parallel program run with one thread per core. For a protocol $P$ on $W$ with
$N$ cores we define
\begin{align*}
  B_{\mathrm{tot}}(P,W,N) &= \text{total bytes moved by coherence messages,}\\
  B_{\mathrm{fwb}}(P,W,N) &= \text{bytes moved by forced writebacks
      (\stM$\to$\stS\ on a remote read),}\\
  T(P,W,N) &= \text{execution time,}\qquad
  \mathrm{AMAT}(P,W,N) \text{ as in \cref{eq:amat-coh},}
\end{align*}
and the forced-writeback fraction of a workload as
\begin{equation}
  F(W,N) = \frac{B_{\mathrm{fwb}}(\mathrm{MESI},W,N)}{B_{\mathrm{tot}}(\mathrm{MESI},W,N)}.
  \label{eq:fwb-fraction}
\end{equation}
$F$ is 30.5\,\% for the trace of \cref{tab:mesi-bill} and between 10\,\% and
45\,\% for the producer--consumer model of \cref{eq:pc-model}. For each $W$ and
$N$ the task is (i)~to measure $F(W,N)$; (ii)~for any protocol $P$ that removes
forced writebacks and keeps both axioms, to measure the change in
$B_{\mathrm{tot}}$, $T$ and AMAT relative to MESI; and (iii)~to find out whether
that change can be predicted from quantities measured under MESI alone.

% ---------------------------------------------------------------
\section{Scope and Assumptions}
\label{sec:prob-scope}
\sectionmap{[Scope [MSI-family write-invalidate] [2 to 16 cores]
  [Consistency model fixed] [OpenMP workloads]]}
\index{MSI}\index{forced writeback}\index{write-back cache}\index{snooping protocol}\index{directory protocol}\index{write-invalidate}\index{false sharing}\index{OpenMP}\index{memory consistency}\index{verification}%

We chose these assumptions so that the forced writeback is the only variable of
interest.
\begin{itemize}[itemsep=1pt]
  \item Protocols: the write-invalidate MSI family (Chapter~5) in snooping and
    directory forms.
  \item Caches: private write-back L1 caches with 64-byte blocks and a shared LLC
    (\cref{sec:cache-private}).
  \item Scale: 2, 4, 8 and 16 cores, where both snooping and directories are
    realistic (\cref{fig:choose-tree}).
  \item Consistency: fixed and not studied (\cref{sec:coh-consistency}).
  \item Workloads: shared-memory C programs parallelised with OpenMP
    (\cref{sec:par-threads}), including kernels built to show each sharing pattern
    of \cref{tab:patterns}.
  \item Out of scope: false sharing (\cref{sec:scale-false}), message passing and
    formal verification.
\end{itemize}

% ---------------------------------------------------------------
\section{Research Questions}
\label{sec:prob-rq}
\sectionmap{[Research questions [RQ1: how big is F] [RQ2: does traffic saved become time saved]
  [RQ3: can it be predicted] [RQ4: when does removing it hurt]]}

The problem splits into four research questions (\cref{tab:rq}).

\begin{table}[htbp]
\centering
\caption{Research questions and where each is answered.}
\label{tab:rq}
\small
\begin{tabular}{@{}l p{8.4cm} l@{}}
\toprule
 & Question & Answered in \\
\midrule
RQ1 & How large is the forced-writeback fraction $F(W,N)$ for realistic
  workloads, and how does it change from 2 to 16 cores? & Part VI \\
RQ2 & When a protocol removes forced writebacks, does the saved traffic turn
  into saved execution time, or is it hidden while the interconnect is lightly
  loaded? & Part VI \\
RQ3 & Can the benefit be predicted from a workload's behaviour under MESI,
  before the other protocol is built? & Parts IV and VI \\
RQ4 & Are there conditions, such as a small LLC, under which deferring the
  writeback costs more than it saves? & Parts IV and VI \\
\bottomrule
\end{tabular}
\end{table}

% ---------------------------------------------------------------
\section{Requirements for Any Solution}
\label{sec:prob-requirements}
\label{sec:prob-bridge}
\sectionmap{[Requirements [Keep completion] [Keep order]
  [Remove the forced write back] [Keep MESI's gains]]}
\index{MESI}\index{MOESI}\index{MESIF}\index{Owned state}\index{forced writeback}\index{directory protocol}\index{token coherence}\index{write-update}\index{coherence axioms}%

A candidate solution in Part~III must meet four requirements, applied in order
in \cref{fig:req-tree}. The first two are the coherence axioms. The third is the
problem itself. The fourth stops a solution from trading one inefficiency for
another. Part~III judges the existing answers against them: the Owned state of
MOESI, which removes the forced writeback directly; write-update protocols, which
avoid invalidation; the Forward state of MESIF, which targets a different cost of
sharing; and directories and token coherence, which change how the cost is paid.

\begin{figure}[htbp]
\centering
\adjustbox{max width=0.85\linewidth}{%
\begin{forest}
  decision tree,
  [Does every deferred write\\still reach memory eventually?\\(completion axiom)
    [Rejected:\\loses writes, el=no]
    [Do all caches still see writes\\to a block in one order?\\(order axiom), el=yes
      [Rejected:\\incoherent, el=no]
      [Does it remove the write back\\on a remote read of an M block?, el=yes
        [Not a solution\\to this problem, el=no]
        [Does it keep MESI's silent\\E-to-M upgrade?, el=yes
          [Weaker candidate, el=no]
          [Candidate for\\Part III, el=yes]]]]]
\end{forest}}
\caption[Decision tree: requirements for a solution]{The four requirements for a
solution, applied in order. Only solutions that reach the bottom-right box are
full candidates in Part~III.}
\label{fig:req-tree}
\end{figure}
